% Tabla generada automaticamente desde la matriz de criterios del codigo.
% No editar a mano; regenerar con genera_tabla_criterios.py (misma carpeta).
{\scriptsize
\begin{longtable}{@{}>{\raggedright\arraybackslash}p{0.238\linewidth}>{\raggedright\arraybackslash}p{0.152\linewidth}>{\raggedright\arraybackslash}p{0.266\linewidth}>{\raggedright\arraybackslash}p{0.295\linewidth}@{}}
\caption{Matriz de criterios estad\'isticos de la aplicaci\'on, generada autom\'aticamente desde el m\'odulo interno de criterios.}
\label{tab:matriz_criterios}\\
\toprule
\textbf{Criterio} & \textbf{Tipo} & \textbf{Valor o umbral} & \textbf{Sustento} \\
\midrule
\endfirsthead
\multicolumn{4}{l}{\emph{Continuaci\'on de la Tabla~\ref{tab:matriz_criterios}}}\\
\toprule
\textbf{Criterio} & \textbf{Tipo} & \textbf{Valor o umbral} & \textbf{Sustento} \\
\midrule
\endhead
\bottomrule
\endlastfoot
Mínimo computacional para modelacion & operativo interno, sin fuente externa & 8 & Parámetro interno documentado \\
Advertencia por serie corta & operativo interno, sin fuente externa & 18 & Parámetro interno documentado \\
Longitud recomendada para modelacion mensual & operativo interno, sin fuente externa & 24 & Criterio metodológico interno \\
Mínimo de iteraciones walk-forward (descriptivo) & operativo interno, sin fuente externa & 6 & Parámetro interno documentado; ninguna fuente fija el valor 6 \\
Primer origen de la validación temporal & pendiente de decisión & N0 = max sobre los candidatos de su mínimo operativo codificado para intentar el ajuste (hoy 6, por Holt amortiguado) & Tashman (2000) y FPP3 §5.10 sustentan el esquema de origen móvil; ninguna fuente determina el primer origen \\
Valor atípico por Z robusto modificado & bibliográfico & |M\_\allowbreak{}i| $>$ 3.5 & NIST/SEMATECH e-Handbook of Statistical Methods, sección 1.3.5.17 \\
Clasificación de alertas por periodo (deduplicadas) & operativo interno, sin fuente externa & patron\_\allowbreak{}calendario / posible\_\allowbreak{}error\_\allowbreak{}datos / posible\_\allowbreak{}cambio\_\allowbreak{}nivel / posible\_\allowbreak{}atipico\_\allowbreak{}aislado & Politica interna; el patron calendario se confirma con C-CAL-001 \\
Subrangos leve/moderado/fuerte de atípicos & operativo interno, sin fuente externa & 3.5-4.5; 4.5-5.0; >5.0 & Sin calibración empírica documentada \\
Durbin-Watson como estadistico descriptivo & bibliográfico & Se reporta el valor en [0, 4]; sin cortes de clasificacion & Durbin y Watson (1951), Biometrika 38(1-2), pp. 159-177 \\
Nivel de significancia Ljung-Box/Jarque-Bera & bibliográfico & alpha = 0.05 & Convencion estadística usual al 5\%; Ljung y Box; Jarque y Bera \\
Heterocedasticidad de los residuos & bibliográfico & Breusch-Pagan sobre [1, t]; se rechaza si p $<$ 0.05 & Breusch y Pagan (1979); statsmodels.stats.diagnostic.het\_\allowbreak{}breuschpagan \\
Media residual igual a cero & bibliográfico & Prueba t de una muestra bilateral; se rechaza si p $<$ 0.05 & Prueba t de una muestra; scipy.stats.ttest\_\allowbreak{}1samp \\
MASE como señal de cautela & bibliográfico & > 1 frente a escala naive in-sample & Hyndman y Koehler (2006) \\
Errores de backtesting inusuales por puntaje z modificado & bibliográfico & |M\_\allowbreak{}i| $>$ 3.5 con M\_\allowbreak{}i = 0.6745 (|e\_\allowbreak{}i| - mediana) / MAD & Iglewicz y Hoaglin (1993); NIST/SEMATECH e-Handbook 1.3.5.17 \\
Construccion del intervalo de prediccion (FPP3 5.5) & bibliográfico & pronostico +/- c * sigma\_\allowbreak{}h, con sigma\_\allowbreak{}h = sqrt(SUM e\textasciicircum{}2 / n) sobre los errores OOS del paso exacto y c el cuantil t con n grados de libertad del nivel nominal & Hyndman y Athanasopoulos (2021), FPP3 5.5: yhat +/- c*sigma\_\allowbreak{}h, con sigma\_\allowbreak{}h estimada de los errores. El multiplicador t es la derivacion exacta cuando sigma\_\allowbreak{}h se estima \\
Condicion de existencia del cuantil de orden al 95 \% & derivación matemática & k = ceil((n+1)(1-alfa)) $\leq$ n, es decir n $\geq$ 1/alfa - 1 = 19 errores al 95 \% & Lei et al. (2018); Angelopoulos y Bates (2023) \\
Cobertura empírica por evaluacion de origen movil & operativo técnico & cada error se contrasta contra el rango construido con los errores ANTERIORES del mismo paso; con n errores se obtienen n-2 contrastes & Christoffersen (1998), International Economic Review 39(4), 841-862, para la evaluacion de una SECUENCIA de intervalos por su sucesion de aciertos y su cobertura incondicional; Tashman (2000) y Hyndman y Athanasopoulos (2021) 5.10 para la evaluacion de origen rodante \\
Ancho relativo de IC95 por horizonte & operativo interno, sin fuente externa & NUEVE cortes: 0.10 operativo; 0.20 corto; 0.30 medio; 0.40 largo; 0.50 extendido cercano; 0.55 extendido; 0.75 exploratorio; mas 0.45 y 0.65 escritos en linea sin constante propia & Politica interna de comunicacion de incertidumbre. Ninguna fuente respalda estos valores \\
Clasificación del intervalo del 95 \% por cobertura empírica & operativo interno, sin fuente externa & ningun corte de cobertura degrada el horizonte; solo degrada que la cobertura no sea calculable & Cobertura observada por origen movil (D-12b-C); Hyndman y Koehler (2006) \\
Seleccion del modelo por RMSE fuera de muestra global & bibliográfico & m* = argmin RMSE\_\allowbreak{}global(m) sobre la muestra comun de pares (objetivo, horizonte) & RMSE: Hyndman y Koehler (2006); FPP3 5.8 (medidas dependientes de escala son las adecuadas dentro de una misma serie). Validacion de origen movil: Tashman (2000); FPP3 5.10. La muestra de evaluacion es el rango entregable h=1..h\_\allowbreak{}max, deducido del contrato de un modelo por trayectoria \\
Continuidad de la evaluacion pese a un horizonte no recomendable (parada retirada) & derivación matemática & el bucle evalua TODOS los horizontes de la rejilla; un horizonte no recomendable no detiene la evaluacion de los siguientes (continue, no break) & Consecuencia tecnica: cada horizonte tiene su propia evidencia fuera de muestra, independiente de la de los demas \\
Horizonte maximo publicado: RETIRADO el prefijo consecutivo & operativo técnico & el maximo publicado es el mayor horizonte que cumple con su propia evidencia; los horizontes validos NO tienen que formar un prefijo & Ninguna fuente exige continuidad; cada horizonte se evalua con su propia muestra de errores \\
Seleccion del horizonte y del modelo entregados & operativo técnico & mayor horizonte permitido menor o igual al solicitado & Regla de entrega; no es un criterio de calidad \\
Salvaguarda por benchmark, retirada como decisora & operativo interno, sin fuente externa & diagnostico: evalua drift y naive y publica el resultado; no sustituye & Decision operativa D-3; el criterio de aceptacion no tenia fuente \\
Horizontes operativos de interfaz & sin efecto decisorio & 1, 3, 6, 12, 18 & Decisión de UX institucional \\
Límite de búsqueda del horizonte & operativo técnico & HORIZONTE\_\allowbreak{}MAXIMO\_\allowbreak{}AUDITORIA = 30 meses, constante & Politica computacional interna \\
Niveles de evidencia fuera de muestra por horizonte & operativo interno, sin fuente externa & 0 ventanas (h $>$ n - N0): no evaluable, por inexistencia del error fuera de muestra; 1 o mas: se entrega, con el numero de ventanas declarado & Derivacion aritmetica: evaluar h con ventana expansiva exige n - N0 - h + 1 $\geq$ 1, luego h $\leq$ n - N0. Por encima no existe ningun error fuera de muestra que medir \\
Parametros de Holt estimados por minimizacion del SSE & bibliográfico & alpha, beta*, l0 y b0 estimados; phi estimado en [0.80, 0.98] en la variante amortiguada; restriccion 0 $<$ beta* $\leq$ alpha & Hyndman y Athanasopoulos (2021), FPP3, secciones 8.1-8.2; Gardner (2006), Exponential smoothing: the state of the art - Part II \\
Elegibilidad de candidatos por estimabilidad (activación por niveles retirada) & derivación matemática & Compiten los 10 candidatos del catalogo; un modelo se excluye solo si n $<$ N\_\allowbreak{}min(m), derivado de su propia formulacion & N\_\allowbreak{}min se deriva por familia, no por una regla unica: forma cerrada (naive, drift) exige solo que existan las observaciones que la formula usa; ajuste por minimizacion (OLS, Huber, Holt) exige n $\geq$ k + 1 contando TODO parametro estimado -en Huber tambien la escala; en Holt tambien alpha, beta*, phi y los dos estados iniciales-; y los modelos sobre variaciones exigen ademas un minimo de variaciones finitas \\
Comparación contra Naive/Drift: constantes retiradas de la logica decisoria & sin efecto decisorio & ningun corte vigente; la lectura viva es C-BEN-002 frente a 1 & Sin fuente para 1,10 ni para 1,25; el inventario RC3 los declaro margen de gracia sin fuente \\
Error relativo frente al benchmark, lectura descriptiva & bibliográfico & no es\_\allowbreak{}benchmark y rRMSE frente a naive $>$ 1; se informa, no decide & Hyndman y Koehler (2006): el error relativo se define frente a UN metodo de referencia y 1 es su punto de equivalencia \\
Detección del salto de cambio de anio & operativo interno, sin fuente externa & $\geq$ 2 transiciones; ratio $>$ 1.5; signo $\geq$ 0.6 & Estadística robusta (medianas); calibrado con el anexo ICOCIV de mayo de 2026 \\
Tratamiento del salto de cambio de año: RETIRADO, no se aplica & operativo interno, sin fuente externa & aplicado = False de forma incondicional; y\_\allowbreak{}ajustado = y\_\allowbreak{}futuro & P0-F, 12-08-2026: ningun componente del tratamiento tiene sustento completo (gamma como estimador del salto futuro, la forma del factor, las puertas de activacion calibradas sobre el propio anexo y el conjunto de validacion) \\
Separacion entre patron detectado y efecto en el horizonte & operativo interno, sin fuente externa & patron\_\allowbreak{}detectado\_\allowbreak{}en\_\allowbreak{}serie, horizonte\_\allowbreak{}cruza\_\allowbreak{}cambio\_\allowbreak{}anio y efecto\_\allowbreak{}en\_\allowbreak{}horizonte\_\allowbreak{}solicitado se informan por separado & Politica interna autorizada el 29-jul-2026 (hallazgo H-10); redaccion completada el 14-08-2026 \\
Estabilidad del error entre ventanas (coeficiente de variación) & operativo interno, sin fuente externa & inestable $>$ 1.0 & Coeficiente de variación de |errores| OOS; parámetro interno \\
\end{longtable}
}
